\documentclass[12pt]{article}
\usepackage{amsmath, amssymb, amsthm}
\usepackage{geometry}
\geometry{a4paper, margin=1in}
\usepackage{hyperref}
\usepackage{enumitem}
\usepackage{graphicx}
\usepackage{float}
\usepackage{listings}
\usepackage{color}
\usepackage{subcaption}

\title{ECE 434: DSP \\ Problem Set 1}
\author{Pranav Chandra N. V}
\begin{document}
	\maketitle

	\begin{enumerate}
		\item Determine the mod-4 and mod-4 reductions of the length 8 signal vector:
			\begin{equation}
				x = [1, 2, -2, 3, 4, -2, -1, 1]^T
			\end{equation}
			\textbf{Solution:}	\\
			Mod-N reduction involves partitioning the vector into sequeces of length N and then summing the sequences element-wise. \\
			Thus, mod-4 reduction of x is given by:
			\begin{equation}
				x_{mod4} = [1+4, 2-2, -2-1, 3+1]^T = [5, 0, -3, 4]^T
			\end{equation}
			Similarly, mod-3 reduction of x is given by:
			\begin{equation}
				x_{mod2} = [1+3-1, 2+4+1, -2-2]^T = [3, 7, -4]^T
			\end{equation}
			Note that for the mod-3, the third sequence has only 2 elements, so we assume the third element is a 0 (pad zeroes where required).

		\item Let $x = [1, 2, 2, 1, 2, 1, 1, 2]$, compute the 4-point DFT using the matrix form definition
		\textbf{Solution:} \\
	\begin{equation}
		\left[
			\begin{matrix}
				1 & 1 & 1 & 1 & 1 & 1 & 1 & 1 \\
				1 & -j&-1&j&1 &-j &-1&  j \\
				1 & -1& 1 &-1&1 & -1& 1 & -1\\
				1 & j & -1& -j&1 & j & -1& -j
			\end{matrix}
		\right]
		\times
		\left[
			\begin{matrix}
			1 \\
			2 \\
			2 \\
			1 \\
			2 \\
			1 \\
			1 \\
			2 \\
			\end{matrix}
		\right]
		=
			\left[
		\begin{matrix}
			12 \\
		0 \\
		0 \\
		0
		\end{matrix}
		\right]
	\end{equation}
	Now, compute the mod-4 reduction of x and then compute the 4-point DFT of the reduced signal. Compare the results with the previous DFT.
	
	\textbf{Solution:} \\
	Mod-4 reduction of x is given by:
	\begin{equation}
		x = [3 , 3, 3, 3]
	\end{equation}
	Now, computing the 4-point DFT of the reduced signal:
\begin{equation}
		\left[
	\begin{matrix}
		1 & 1 & 1 & 1 \\
		1 & -j & -1 & j \\
		1 & -1 & 1 & -1 \\
		1 & j & -1 & -j 
	\end{matrix}
	\right]
	\times
		\left[
	\begin{matrix}
		3 \\
		3 \\
		3 \\
		3 
	\end{matrix}
	\right]
	 = 
	 	\left[
	 \begin{matrix}
		 12 \\ 
		 0 \\
		 0 \\
		 0
	 \end{matrix}
	 \right]
\end{equation}
	
This question basically shows that to compute reduced points DFT, taking the mod-N version of the main sequence provides you with a more efficient way of computing the DFT.

	\item Comute the 8 point FFT of a length 8 signal which are8 samples of the signal:
		\begin{equation}
			x[n] = 4\cos\left(\frac{xn}{2}\right) + \cos(\pi n)
		\end{equation}
		Discuss if the FFT accurately represents the entire sepctrum of $n]$.

		\textbf{Solution:} \\
		The 8 points sequence to be taken is given by:
		\begin{equation}
			x = [ 0, 0, 16, 0, 8, 0, 16, 0]
		\end{equation}
		Now, computing the 8-point DFT of the signal:
		
	\item Find the DTFT of $y_1[n] = \begin{cases} 1 \hspace{5mm} \text{if } -N \leq n \leq N \\ 0 \hspace{5mm} \text{otherwise} \end{cases}$.

		\textbf{Solution:} \\
	The DTFT of a signal is given by:


	\begin{equation}
		Y_1(e^{j\omega}) = \sum_{n=-\infty}^{\infty} y_1[n] e^{-j\omega n}
	\end{equation}

	\item Let $X(k)$ be a 14-oitn DFT of a length 14 real sequence $x(n)$. The first 8 sampels of X(k) are given by:
		\begin{align}
			X(0) = 12 \\
			X(1) = -1 + j3 \\
			X(2) = 3 + j4 \\
			X(3) = 1 - j5 \\
			X(4) = -2 + j2 \\
			X(5) = 6 + j3 \\
			X(6) = -2 - j3 \\
			X(7) = 10
		\end{align}
		\begin{enumerate}
			\item Find the remaining 6 values of X(k). \\
				\textbf{Solution:} \\
				Using the property of DFT for real sequences, we know that:
				\begin{equation}
					X(N-k) = X^*(k)
				\end{equation}
				where N is the length of the sequence. \\
				Thus, we can compute the remaining values as follows:
				\begin{align}
					X(8) = X^*(6) = -2 + j3 \\
					X(9) = X^*(5) = 6 - j3 \\
					X(10) = X^*(4) = -2 - j2 \\
					X(11) = X^*(3) = 1 + j5 \\
					X(12) = X^*(2) = 3 - j4 \\
					X(13) = X^*(1) = -1 - j3
				\end{align}
			\item Find $x(0)$ without computing the IDFT of X(k). \\
				\textbf{Solution:} \\
				Using the property of DFT, we know that:
				\begin{equation}
					x(0) = \frac{1}{N} \sum_{k=0}^{N-1} X(k)
				\end{equation}
				where N is the length of the sequence. \\
				Thus, we can compute $x(0)$ as follows:
			\item Find $x(7)$ without computing the IDFT of X(k). \\
				\textbf{Solution:} \\
				Using the property of DFT, we know that:
				\begin{equation}
					x(n) = \frac{1}{N} \sum_{k=0}^{N-1} X(k) e^{j\frac{2\pi}{N}kn}
				\end{equation}
				where N is the length of the sequence. \\
				Thus, we can compute $x(7)$ as follows:
		\end{enumerate}
		
	\item A length 8 sequence $x(n) = \{-4, 5, 2, -3, 0, -2, 3, 4\}$ has an 8-point DFT $X(k)$. A new sequence is defined as:
		\begin{equation}
			Y(k) = W_4^{3k} X(k)
		\end{equation}
		Find y(n) without computing the IDFT of Y(k). \\
		\textbf{Solution:} \\	
		We know that for a DFT, multiplying by a twiddle factor corresponds to a circular shift in the time domain. \\
		Thus, we can compute y(n) as follows:
		\begin{equation}
			y(n) = x((n-3) \mod 8)
		\end{equation}
		This process exists because multiplying by $W_N^{kn_0}$ in the frequency domain corresponds to a circular shift by $n_0$ in the time domain. However, when $N$ is different from the length of the sequence, the circular shift is scaled accordingly. In this case, since we are multiplying by $W_4^{3k}$ instead of $W_8^{3k}$, the effective shift in the time domain is scaled by a factor of 2 (since 8/4 = 2). Therefore, the circular shift is by 3*2 = 6 positions.

	\item Let $X(z) = Z\{x(n\}$, where $x(n) = 0.4^nu(n)$:
		\begin{enumerate}
			\item Find $g(n) = Z^{-1}\{X(z^2)\}$ without computing $X(z)$.
			\item Find $g(n) = Z^{-1}\{(1 + Z^{-1})X(z^2)\}$ without computing $X(z)$.
		\end{enumerate}
	\item Let $X(z) = Z\{x(n)\}$, $x(n)$ is a 12 length sequence, and $X_0(k)$ is the 9 point DFT of $x(n)$. Find the inverse DFT of $X_0(k)$ in terms of $x(n)$ without computing $X_0(k)$.
	\item Let $X(e^{j\omega}) = F.T.\{1, 1, ..., 1\}$ from 0 to 7. $X_1(k)$ is obtained by sampling $X(e^{j\omega})$ at uniform intervals of $\frac{\pi}{5}$. 
	\end{enumerate}

\end{document}
